% Adopted Lawvere continuation; working spec fragment preserved separately.
\begin{annexitem}{annex:p10.lawvere-bvm.base}{thm:p10-lawvere-bvm-base}{The base and bounded well-ordering}
\begin{proof}
If the topos were degenerate, every object and every endomap would be
uniquely isomorphic to the terminal one. The endomap forgetful functor
would be an equivalence, contrary to the source's \((\omega)\).
Classical complements and irreducibility give, for every subterminal
\(U\), either \(U=1\) or \(\neg U=1\), hence \(U=0\).
The already proved support/classicality generator theorem then makes
\(1\) a generator. The proved \((\Pi)+(\exists)\) equivalence
makes every epi split.

Here is a bounded proof of well-ordering, needed below rather than an
unproved appeal to a set-theoretic universe. For an object \(I\), let
\(\mathcal P_+I\) be the subobject of nonempty predicates. Membership
\(\{(A,x):x\in A\}\to\mathcal P_+I\) is epi, by the definition
of nonempty and the established image calculus. Split it to obtain
\(c(A)\in A\). Everything that follows is a bounded internal
classical argument, using subobjects of \(I\) and \(I\times I\).
Call a well-order on a subset \(D\subseteq I\) admissible if
\[
       x=c(I\setminus\{y:y<x\})\quad(x\in D).
\]
The collection of admissible domain/relation pairs is a subobject of
\(\mathcal P(I)\times\mathcal P(I\times I)\), so its union
domain and union relation exist; singleton and empty domains are not
lost by encoding only a strict relation.

Any two admissible orders have matching initial segments and one is
an initial segment of the other. To verify this without ordinal
comparison, take the union \(C\) of their common initial segments
on which their relations agree. The union is again such a segment:
two constituent segments are nested in each well-order, and comparison
with any element is unchanged. If both orders have elements outside
\(C\), their least outside elements are both \(c(I\setminus C)\).
Adjoining that same element enlarges the common initial segment, a
contradiction. Thus one order ends at \(C\).

Let \(D\) and \(<\) be the union of all admissible orders.
Comparability just proved makes this a linear order with each constituent
an initial segment. For a nonempty \(A\subseteq D\), take
\(x\in A\) and a constituent containing it. The least element of
\(A\) in that constituent is least in all of \(A\): any smaller
element belongs to the same initial segment. This proves well-ordering;
the predecessors of an element are unchanged in larger constituents,
so the union is admissible. If \(D\ne I\), append
\(c(I\setminus D)\) to this admissible order. The resulting
admissible relation belongs to the collection whose union was \(D\),
a contradiction. Thus \(D=I\). The empty object has its empty order.
All unions here are images of bounded membership relations; there is
no union over external ordinals. The internal proof yields a nonempty
object of well-orders, and epi splitting chooses an actual order if
needed. No recursion-lifting or free-algebra theorem was used.
\end{proof}
\end{annexitem}

\begin{annexitem}{annex:p10.lawvere-bvm.presentation}{thm:p10-lawvere-bvm-presentation}{Complete truth algebras and effective matrix presentations}
\begin{proof}
Write \(T=T_{\mathcal E}\). Adjunction and classification give,
naturally in every \(I\in\mathcal S\),
\[
 \mathcal S(I,p_*T)\cong\mathcal E(p^*I,T)
              \cong\operatorname{Sub}_{\mathcal E}(p^*I).
\tag{1}
\]
Finite products preserved by \(p_*\) transport Boolean operations
and their equations to \(B=p_*T\). For a family \(a:I\to B\),
transpose to \(\bar a:p^*I\to T\); its join is the subterminal
image of its classified subobject under \(p^*I\to1\). More
generally, with parameters \(K\) and a family over \(K\), use
the image along \(p^*I\to p^*K\). Image minimality and pullback
stability prove the upper-bound and least-upper-bound assertions in
every parameter context. Applied to generic membership in
\(\mathcal P_{\mathcal S}B\times B\), this gives an actual
map \(\bigvee:\mathcal P_{\mathcal S}B\to B\), not merely
external joins of global elements. The construction commutes with
reindexing by the image pullback law. Boolean complementation supplies
meets. In particular finite meets distribute over these joins: meet
with \(a\) is left adjoint to \(\neg a\vee(-)\), and therefore
preserves least upper bounds.

We construct an effective presentation of every \(X\in\mathcal E\).
Let \(\partial X\subseteq\mathcal P_{\mathcal E}X\) consist of
subsingletons. Its universal member is the graph of a partial element:
its projection to \(\partial X\) is monic, since the members are
subsingletons. Put \(I_X=p_*\partial X\). Pull this monic
domain back along the counit \(p^*I_X\to\partial X\), obtaining
\(D_X\hookrightarrow p^*I_X\) and evaluation \(d_X:D_X\to X\).
This map is epi. Indeed any arrow \(U\to X\) with \(U\)
subterminal determines a subsingleton subobject of \(X\), hence a
global partial element \(1\to\partial X\). By adjunction it
occurs in the displayed family and its domain is exactly \(U\).
If two maps agree after \(d_X\), they therefore agree after every
subterminal arrow into \(X\); the proved source generator theorem
makes them equal. This is epi cancellation, not an assumption about
global points of \(\mathcal E\).

Under (1), \(D_X\) has weights \(b_i\), and its kernel pair
has a predicate \(e_{ii'}\) on \(p^*(I_X\times I_X)\).
It satisfies precisely the object equations in the matrix definition.
The established regularity and effective quotient theorem gives
\(X=D_X/e\). Conversely, from any \((I,b,e)\) use (1) to form
\(D_b\hookrightarrow p^*I\) and the equivalence relation classified
by \(e\) on \(D_b\). Reflexivity, symmetry and transitivity follow
from the displayed equations, so its effective quotient \(Q(I,b,e)\)
exists. The discrete presentation of \(p^*I\) has \(b_i=1\)
and \(e_{ii'}\) equal to the truth of \(i=i'\) in the base;
\(p^*\) preserves this diagonal and its Boolean complement.

Here are the actual maps, not just the object assignment. A matrix
\(r\) classifies a relation \(R\subseteq D_b\times D_c\).
Its two saturation laws say that membership is invariant under the two
kernel pairs. Its image \(G\subseteq Q(I,b,e)\times Q(J,c,f)\)
pulls back to exactly \(R\). For example, a locally chosen lift of
each quotient coordinate changes that lift only by \(e\) or \(f\);
saturation gives both inclusions. This local statement is image descent
along the product of the two effective quotient epis, established by
their kernel pairs. Totality makes \(G\to Q(I,b,e)\) epi.
Single-valuedness makes that same map mono: two output lifts with a
common input have \(f\)-equivalent values. It is therefore an
isomorphism, and \(G\) is the graph of a unique map. Conversely
pulling a map's graph back along these two covers gives all five
matrix laws. These operations are inverse; equality is tested after
the product epi. Pullback of a composite graph gives
\(\bigvee_j r_{ij}\wedge s_{jk}\), because the middle quotient
is covered by \(D_c\). Pullback of an identity gives \(e\).
This proves the functor laws, associativity and both unit laws.

Thus \(Q:\operatorname{Mat}(B)\to\mathcal E\) is fully faithful.
The canonical partial-element covers give a functor \(M\) in the
other direction: on a map use the pulled-back graph just described.
The counit \(QM X\to X\) is the unique map whose composite with
\(D_X\to QM X\) is \(d_X\); effectiveness makes it an
isomorphism. The unit \(A\to MQ A\) is the matrix of the relation
``these two local representatives have the same quotient value.''
Its inverse is the converse relation. Their composites are the two
kernel-pair matrices, because each quotient cover is epi. Naturality
is graph composition. Applying either comparison to a representative
does not change its quotient value, so both triangle identities hold
after the respective epi covers and hence hold outright. On discrete
presentations these maps give the stated comparison with the base.
\end{proof}
\end{annexitem}

\begin{annexitem}{annex:p10.lawvere-bvm.realization}{thm:p10-lawvere-bvm-realization}{Internal Boolean sheaves and the explicit choice function}
\begin{proof}
Regard \(B\) as an internal thin category. For a sieve \(S\)
on \(b\), put \(jS=\{c\le b:c\le\bigvee S\}\). This is
idempotent, fixes the maximal sieve and preserves intersections:
\(\bigvee(S\cap R)=(\bigvee S)\wedge(\bigvee R)\).
For the nontrivial inclusion, distribute both joins; every
\(s\wedge r\) lies in the intersection of the downward-closed
sieves. Restriction commutes with \(j\), by the same distribution.
Thus this is an actual topology on the internal sieve classifier.
The already constructed internal-presheaf topos and associated-sheaf
reflector give \(\operatorname{Sh}_{\mathcal S}(B)\) without an
external family of products or a Grothendieck-universe assumption.

Constant families have the internal-sections right adjoint by the
proved internal restriction/extension adjunction. Therefore
\(q^*=a\Delta\) has right adjoint \(\Gamma i\), and is left
exact. The closed-sieve classifier at \(b\) is the object
\(\{c:c\le b\}\), with restriction \(c\mapsto c\wedge b'\).
It is Boolean, its complement being relative complement
\(b\wedge\neg c\). Evaluation at the top gives
\(\Gamma T=B\), with its original operations. The same calculation
with a parameter \(I\) says that subobjects of \(q^*I\) are
exactly predicates \(I\to B\): a closed subobject of the constant
presheaf is determined by, for each \(i\), the principal downset
generated by the join of its stages. Conversely these downsets define
a closed subobject. Consequently \(q_*T=B\) internally.

The representable \(y b\) is a subterminal sheaf. A covering family
of stages has a map to \(y b\) exactly when every stage is below
\(b\); if their join is below \(b\), this condition holds and
the unique map glues. For a sheaf \(F\), let \(I\in\mathcal S\)
be the total object of its internal stage-elements \((b,x)\).
The domains with weights \(b\) form a subobject \(D\subseteq q^*I\)
and the internal Yoneda evaluations give \(D\to F\). Its image is
all of \(F\): at every generalized stage \((b,x)\), the element
\(x\) itself is in that image. This is an internal, parameterized
stage test, not a test on external points. Pulling back graphs to
such covers and taking effective quotients gives exactly the matrix
equivalence and its two comparison maps proved in the preceding annex.
That argument required only the existence of these covers, not source
Skolem or choice. Hence it applies here without circularity.

We prove choice by a calculation valid for either matrix presentation.
Let an epi \(Y=(J,c,f)\to X=(I,b,e)\) have matrix
\(v_{ji}\). Pulling back its image along the target cover gives
\[
                   \bigvee_j v_{ji}=b_i.                 \tag{2}
\]
Conversely (2) makes the image all of \(X\), hence makes the map
epi. Choose a well-order of \(J\) in \(\mathcal S\), using the
preceding bounded construction, and define
\[
 t_{ij}=v_{ji}\wedge\neg\bigvee_{k<j}v_{ki},\qquad
 r_{ij'}=\bigvee_j t_{ij}\wedge f_{jj'}.                \tag{3}
\]
For fixed \(i\) the \(t_{ij}\)'s are disjoint and their join
is \(b_i\). Here is the latter assertion for an arbitrary internal
well-order, not just a finite enumeration. Put
\(a=b_i\wedge\neg\bigvee_j t_{ij}\). By well-founded induction
on \(j\), \(a\wedge v_{ji}=0\): if this holds below \(j\),
distribution gives \(a\wedge\bigvee_{k<j}v_{ki}=0\), hence
\(a\wedge v_{ji}\le t_{ij}\), which is disjoint from \(a\).
Taking the join and using (2) gives \(a=0\). Well-founded induction
itself follows by taking the least member of a nonempty counterexample
subobject in \(\mathcal S\).

The matrix \(r\) is supported and total. Its output saturation is
transitivity of \(f\); its single-valuedness follows since distinct
\(t_{ij}\)'s are disjoint, and for a common \(j\),
\(f_{jj'}\wedge f_{jj''}\le f_{j'j''}\). For input saturation,
the saturation of \(v\) implies
\(e_{ii'}\wedge v_{ji}=e_{ii'}\wedge v_{ji'}\) for every
\(j\). Distribution and relative Boolean complements then give
the same equality for \(t_{ij}\), and hence for \(r\).
Finally
\[
 (vr)_{ii'}=\bigvee_j t_{ij}\wedge v_{ji'}
           =\bigvee_j t_{ij}\wedge e_{ii'}=e_{ii'}.
\]
The first equality uses saturation of \(v\) under \(f\).
For the second, \(t_{ij}\le v_{ji}\) and single-valuedness
give one inclusion; saturation under \(e\) gives the other.
The last uses \(\bigvee_j t_{ij}=b_i\) and \(e_{ii'}\le b_i\).
Thus (3) is an actual section. All indices and induction are in
\(\mathcal S\); no representatives of sheaf elements were chosen.
Every epi splits, so the proved choice equivalence gives
\((\Pi)\) and \((\exists)\). Classicality was already obtained
from the explicit classifier, establishing the exact source package.
If \(B\) is the one-element algebra, all matrix predicates coincide
and every object and map is uniquely isomorphic to the terminal one;
the sheaf construction agrees. This verifies the degenerate boundary.
\end{proof}
\end{annexitem}

\begin{annexitem}{annex:p10.lawvere-bvm.correspondence}{thm:p10-lawvere-bvm-correspondence}{The morphism correspondence, its adjunctions and its 2-cells}
\begin{proof}
For a complete homomorphism \(\alpha:B\to C\), applying
\(\alpha\) to every entry preserves all object and map equations.
It preserves identities and composition by preservation of joins and
finite meets. This defines \(T_\alpha:\operatorname{Mat}(B)\to
\operatorname{Mat}(C)\), including its actual action on maps.
It fixes the discrete base presentations, since it fixes \(0,1\)
and every Boolean predicate of the base takes these two values.

This functor is left exact. The terminal presentation is a singleton
of weight \(1\). The product of \((I,b,e)\) and \((J,c,f)\)
has indices \(I\times J\), weights \(b_i\wedge c_j\), and
equality \(e_{ii'}\wedge f_{jj'}\); its projections have the
graphs of the coordinate maps. Given parallel matrices \(r,s:A\to D\),
their equalizer has the same indices as \(A\), weights
\(a_i=\bigvee_j r_{ij}\wedge s_{ij}\), and equality
\(a_i\wedge e_{ii'}\wedge a_{i'}\). This is the predicate
that their two quotient values agree: output saturation incorporates
all pairs of equivalent output representatives. The universal property
follows by the graph criterion of the presentation proof. These formulas
are preserved by \(\alpha\), proving left exactness.

The right adjoint is not assumed from an adjoint-functor theorem.
For a \(C\)-sheaf \(Y\), define
\[
             R_\alpha Y(b)=Y(\alpha b).
\]
Its restriction maps are those of \(Y\). A covering family of
\(b\) maps to one of \(\alpha b\), because \(\alpha\)
preserves joins and meets, so a matching family has a unique glue.
Thus \(R_\alpha Y\) is a \(B\)-sheaf. A matrix presentation
\(A=(I,b,e)\) of a sheaf is the effective quotient of the
base-indexed family of representables \(y b_i\), with overlaps
\(y e_{ii'}\). This means the subobject of the discrete family
with these weights, not an assumed external coproduct. A map from
its quotient into any sheaf is precisely an internal family
\(s_i\) of sections at \(b_i\) agreeing after restriction to
\(e_{ii'}\). Existence is effective epi descent; uniqueness is
epi cancellation. Accordingly
\[
 \operatorname{Hom}(T_\alpha A,Y)
 \cong\{(s_i):s_i\in Y(\alpha b_i),\ s_i|_{\alpha e_{ii'}}
                           =s_{i'}|_{\alpha e_{ii'}}\}
 \cong\operatorname{Hom}(A,R_\alpha Y).                 \tag{4}
\]
These are inverse natural correspondences, at every parameter base.
The unit sends each local representative to its same section in the
middle family in (4); the counit evaluates that family in \(Y\).
Their composites on each local representative are the identity.
The presentation covers are epi, so both triangle identities hold.
This supplies the right adjoint, unit and counit and hence an actual
geometric inverse image over \(\mathcal S\).

Conversely let \(L:\mathcal E\to\mathcal F\) be a specified
1-cell over the base. Transport a subobject of \(p^*I\) by \(L\)
and its base isomorphism to a subobject of \(q^*I\).
By (1) this gives natural maps
\(\mathcal S(I,B_{\mathcal E})\to
\mathcal S(I,B_{\mathcal F})\), hence an actual
\(\alpha_L:B_{\mathcal E}\to B_{\mathcal F}\) by Yoneda in
\(\mathcal S\). This preserves finite meets and Boolean
complements: \(L\) preserves finite limits, disjoint coproducts,
and their covering images. It preserves base-indexed joins as well.
Such a join is the image of a subobject of \(p^*I\) under its
projection (with arbitrary parameters). A finite-limit-preserving
left adjoint preserves effective epis and the epi/mono image
factorization, since an epi is the coequalizer of its kernel pair.
It therefore preserves exactly that image. This proves completeness,
not just preservation of finite Boolean operations.

Apply \(L\) to the effective presentation \(D_b/R_e\).
It takes \(D_b\) to \(D_{\alpha_L b}\), takes \(R_e\)
to \(R_{\alpha_L e}\), and preserves their effective quotient.
On maps the same assertion follows by transporting the graph and its
image. Thus there is a natural isomorphism
\(L\cong T_{\alpha_L}\), under the two matrix equivalences.
It is the identity on local quotient values and is compatible with
the specified base isomorphism. For \(T_\alpha\), transport of
a weighted subterminal gives \(\alpha(b)\), so
\(\alpha_{T_\alpha}=\alpha\). For the model realized from \(B\),
the algebra comparison is the explicit classifier evaluation
\(q_*T\cong B\) proved above. For an original model, the
model comparison is the quotient equivalence \(QM\cong1\),
with its proved inverse and triangles. These are the object and
1-cell comparisons of the claimed correspondence.

It remains essential to check 2-cells rather than declare a category
equivalence from these objects. Let \(\tau:L\Rightarrow L'\)
be compatible with the fixed base comparison. On every subterminal
\(U_b\to1\), naturality gives a map
\(U_{\alpha b}\to U_{\beta b}\) over \(1\); hence
\(\alpha b\le\beta b\). Apply this also to \(\neg b\).
Complement preservation reverses the resulting inequality, so
\(\alpha b=\beta b\). This holds internally by the same argument
on the universal weighted subterminal with base index \(B\).
Thus \(\alpha=\beta\), not merely equality on global truth values.

On a weighted subobject \(D\to p^*I\), naturality and base
compatibility force \(\tau_D\) to be the unique inclusion between
the equal transported subobjects. On its effective quotient, naturality
forces \(\tau\) to be the quotient comparison. Since every object
has such a presentation, there is at most one compatible 2-cell; it is
an isomorphism when the homomorphisms agree. Conversely the two
presentation comparisons for equal homomorphisms produce that unique
2-cell. Thus every hom-category is equivalent to the discrete set of
complete homomorphisms. The right adjoint mate is determined by (4)
and introduces no further independent cells.

Finally \(T_\beta T_\alpha=T_{\beta\alpha}\) and
\(T_1=1\) on matrices, entry by entry. For transported sheaf or
model presentations the composite comparison is the map preserving
the same local quotient value. Both sides of each associativity
pentagon and unit triangle preserve that value, so epi cancellation
makes them equal. The adjunction comparisons agree by (4) and its
triangle identities. This proves pseudofunctor coherence and the
biequivalence in the declared inverse-image direction. Reversing
1-cell direction gives the conventional geometric opposite. None
of this assumes that these left adjoints preserve powers: that stronger
logical-morphism property is not part of this morphism regime.
\end{proof}
\end{annexitem}
